\documentclass[11pt,a4paper]{article}

\usepackage[utf8]{inputenc}
\usepackage[T1]{fontenc}
\usepackage[danish]{babel}
\usepackage{lmodern}
\usepackage[a4paper,margin=2.5cm]{geometry}
\usepackage{enumitem}
\usepackage{xcolor}
\usepackage{amssymb}
\usepackage{hyperref}

\hypersetup{
  colorlinks=true,
  linkcolor=blue!55!black,
  urlcolor=blue!55!black,
  pdftitle={Arkitektur og design af DTU Course API}
}
\setlist{nosep,leftmargin=*}
\setlength{\parindent}{0pt}
\setlength{\parskip}{0.65em}
\newcommand{\code}[1]{\texttt{#1}}

\title{DTU Course API\\[0.3em]
\large Arkitektur, design og søgeforløb}
\author{}
\date{August 2026}

\begin{document}

\maketitle

\begin{abstract}
Projektet er en webbaseret studievejleder, som gør officielle DTU-data
tilgængelige gennem både en chat og et REST-API. Systemet kan søge efter
kurser, slå et kursusnummer op og forklare studieplaner og specialiseringer.
Denne rapport beskriver de centrale arkitekturvalg og følger en forespørgsel
fra brugerens besked til det svar, der vises i browseren. Fokus er dermed på
systemets ansvar og dataflow frem for en gennemgang af hver enkelt kildefil.
\end{abstract}

\tableofcontents

\section{Formål og afgrænsning}

Systemets formål er at give studerende en enkel indgang til DTU's
kursuskatalog og studieordninger. Brugeren behøver ikke kende API-parametre
eller den præcise struktur på DTU's sider, men kan eksempelvis skrive:

\begin{itemize}
  \item ``Giv mig alle kurser i cyber''.
  \item ``Hvad handler kursus 02150 om?''
  \item ``Computer Science study plan''.
  \item ``Hvilke specialiseringer findes på Computer Science and Engineering?''
\end{itemize}

Applikationen er ikke et generelt søgesystem til hele internettet. Svar om
kurser, studieplaner og specialiseringer skal baseres på de officielle data,
som allerede er importeret til projektets database. Links i svaret fører
brugeren tilbage til DTU som autoritativ kilde.

\section{Overordnet arkitektur}

Løsningen er opbygget som en lagdelt webapplikation. FastAPI er både
webserver, REST-API og vært for et read-only MCP-interface. PostgreSQL lagrer
de normaliserede DTU-data og understøtter både fuldtekstsøgning og
vektorsøgning gennem pgvector.

\begin{figure}[ht]
\centering
\fbox{\begin{minipage}{0.92\textwidth}
\centering
\textbf{Browser: HTML, CSS og JavaScript}\\
$\downarrow$ \code{POST /api/chat}\\[0.25em]
\textbf{FastAPI: validering og chat-route}\\
$\downarrow$\\[0.25em]
\textbf{Recommendation service: sprog, intent og orkestrering}\\
$\swarrow$ \hspace{4em} $\downarrow$ \hspace{4em} $\searrow$\\[0.25em]
\textbf{Søgeservice} \hspace{2em}
\textbf{Groq + MCP} \hspace{2em}
\textbf{Studieplanslogik}\\
$\searrow$ \hspace{4em} $\downarrow$ \hspace{4em} $\swarrow$\\[0.25em]
\textbf{PostgreSQL/pgvector med importerede DTU-data}
\end{minipage}}
\caption{De vigtigste komponenter i behandling af en chatbesked.}
\end{figure}

Ansvaret er fordelt således:

\begin{description}[style=nextline]
  \item[Præsentationslaget]
  Filerne i \path{app/web/} viser chatten, sender brugerbeskeder til API'et og
  renderer almindelig tekst samt strukturerede kursus-, studieplans- og
  specialiseringskort.

  \item[API-laget]
  Routes i \path{app/api/routes/} validerer input med Pydantic og opretter en
  SQLAlchemy-session pr. request. Den offentlige chat bruger
  \code{/api/chat}; det separate kursus-API under \code{/api/v1} kræver en
  \code{X-API-Key}.

  \item[Servicelaget]
  Filerne i \path{app/services/} klassificerer brugerens hensigt, finder
  studieprogrammer, udfører søgninger og sammensætter svaret. Routes holdes
  derfor forholdsvis tynde.

  \item[AI- og værktøjslaget]
  Groq anvendes til semantisk intent-klassifikation og til formulering af
  svar. Ved faktuelle spørgsmål får modellen kun databaseadgang gennem fire
  MCP-værktøjer: \code{get\_course}, \code{search\_courses},
  \code{get\_study\_plan} og \code{get\_specializations}. OpenAI's embedding-
  API bruges særskilt til at omdanne en søgetekst til en vektor.

  \item[Datalaget]
  SQLAlchemy-modeller repræsenterer kurser, oversættelser, studieprogrammer,
  sektioner, krav og specialiseringer. Alembic styrer databaseskemaet.
\end{description}

Denne opdeling isolerer brugergrænsefladen fra søgealgoritmen og databasen.
Det betyder eksempelvis, at den semantiske søgning kan ændres uden at ændre
chatformatet.

\section{Data og importdesign}

Data indlæses på forhånd af modulerne i \path{importer/}. Kursusdata hentes
fra gemte DTU XML-data, mens studieplaner og specialiseringer parses fra deres
respektive DTU-sider. Importerne bruger indholdshashes, så uændrede data ikke
skal behandles igen, og de registrerer fejl uden at gøre den offentlige chat
afhængig af et live-kald til DTU.

De vigtigste relationer er:

\begin{itemize}
  \item Et \code{Course}-objekt identificeres af kursusnummer og studieår og
  indeholder blandt andet ECTS, niveau, skemagruppe og kilde-URL.
  \item Hvert kursus har separate \code{CourseTranslation}-rækker for dansk
  og engelsk tekst. Her ligger også søgevektor og embedding.
  \item Et \code{StudyProgram} består af ordnede sektioner, kurser og krav.
  Krav kan blandt andet udtrykke et ECTS-minimum, et bestemt antal kurser
  eller at alle kurser skal tages.
  \item En \code{StudySpecialization} tilhører et studieprogram og har sine
  egne kurser og krav. De er adskilt fra programmets generelle krav, så en
  valgfri specialisering ikke fejlagtigt præsenteres som obligatorisk for alle.
\end{itemize}

Den normaliserede model gør det muligt at returnere strukturerede resultater
i stedet for kun at lade en sprogmodel genfortælle en lang HTML-side.

\section{Behandling af en chatbesked}

Når brugeren sender en besked, sker følgende:

\begin{enumerate}
  \item Browseren gemmer beskeden lokalt og sender højst de seneste 12
  \emph{brugerbeskeder}. Tidligere assistentsvar sendes ikke tilbage.
  \item Pydantic validerer requesten. Hver besked skal indeholde 1--800 tegn,
  og studieåret skal have formatet \code{YYYY-YYYY}.
  \item Backend registrerer sproget i den seneste besked og klassificerer
  også intent ud fra netop denne besked.
  \item Hele rækken af brugerbeskeder samles som kontekst. Den bruges til at
  genfinde oplysninger, som er udeladt i en opfølgning, eksempelvis navnet på
  et studieprogram.
  \item Det valgte søgeflow henter fakta fra databasen, eventuelt gennem MCP.
  \item API'et returnerer et \code{ChatResponse} med tekstsvaret og eventuelle
  strukturerede anbefalinger, studieplaner eller specialiseringer.
  \item Frontend viser svaret og vælger danske eller engelske overskrifter og
  links ud fra feltet \code{responseLanguage}.
\end{enumerate}

At den seneste besked styrer intent er et vigtigt designvalg. Hvis brugeren
først søger efter cyberkurser og bagefter skriver ``Find kurser i machine
learning'', behandles anden besked som en ny kursussøgning. Resultatet fra den
første søgning må ikke overtage klassifikationen af den næste.

\section{Søgetyper og deres dataflow}

\subsection{Emnebaseret kursussøgning}

En formulering som ``Find kurser i machine learning'' klassificeres som en
kursusanbefaling. Kendte emner genkendes deterministisk. Hvis den regelbaserede
klassifikation ikke er sikker, kan Groq i stedet producere en valideret
\code{SemanticQueryPlan}; modellen vælger kun domæne, operation og emner og
leverer ikke selve kursusfakta.

I det normale AI-flow sendes samtalekonteksten til Groq. Modellen kalder
\code{search\_courses} gennem MCP, og værktøjet udfører den faktiske søgning i
PostgreSQL. Søgeemnet normaliseres til en kort engelsk frase, mens
\code{search\_language} bestemmer sproget på titler og beskrivelser.

Hvis Groq eller MCP ikke svarer, falder systemet tilbage til sin lokale
fortolkning af emne, niveau, ECTS, sprog og periode. Her returneres op til fem
kurser. Niveau- og ECTS-filteret lempes trinvist, hvis den første søgning ikke
giver resultater.

\subsection{Hybrid tekst- og vektorsøgning}

Selve \code{search\_courses}-servicen kombinerer to forskellige signaler:

\begin{enumerate}
  \item \textbf{Leksikalsk søgning:} PostgreSQL omdanner forespørgslen til en
  full-text query og rangerer matches i de danske eller engelske kursustekster.
  Et GIN-indeks gør dette effektivt.
  \item \textbf{Semantisk søgning:} OpenAI-modellen
  \code{text-embedding-3-small} omdanner forespørgslen til en vektor med 1536
  dimensioner. pgvector beregner cosine-afstanden til de embeddings, som blev
  genereret for kursusteksterne under en separat backfill-proces.
  \item \textbf{Fletning:} Resultatlisterne kombineres med Reciprocal Rank
  Fusion (RRF). Metoden bruger placeringen i hver liste frem for at sammenligne
  to inkompatible rå scorer direkte.
\end{enumerate}

Den semantiske del køres kun med PostgreSQL, aktiveret semantisk søgning og en
konfigureret embedding-nøgle. Fejler embedding-kaldet eller vektorsøgningen,
bevares den leksikalske resultatliste. En søgning kan derfor stadig fungere,
selv om OpenAI midlertidigt er utilgængelig.

\subsection{Forespørgsler efter alle match}

Ord som ``alle'', ``every'' eller ``complete list'' aktiverer et særligt flow.
Den semantiske plan kan udtrække et eller flere emner. Systemet søger hvert
emne med en høj intern grænse, samler resultaterne, fjerner dubletter på
kursusnummer og sorterer den endelige liste stigende. Dette omgår MCP-
værktøjets normale maksimum på 20 resultater og gør ``alle kurser i cyber''
forskellig fra en kort anbefaling.

\subsection{Direkte opslag på kursusnummer}

Et femcifret tal giver højeste sikkerhed og klassificeres som et direkte
kursusopslag. Backend kontrollerer først, at kurset findes i det valgte
studieår. Groq får derefter adgang til det konkrete kursus gennem
\code{get\_course} og kan besvare spørgsmål om eksempelvis indhold,
forudsætninger, eksamen eller undervisere. Hvis nummeret ikke findes,
returneres en præcis ``course not found''-besked i stedet for en bred søgning.

\subsection{Studieplaner}

Ord som ``studieplan'', ``study plan'', ``obligatorisk'' og ``ECTS'' fører til
studieplansflowet. Programmet matches mod navn, importerede aliaser, slug og
uddannelsesniveau. Mindre stavefejl kan håndteres med fuzzy matching; hvis det
stadig er uklart, kan en semantisk resolver vælge blandt de faktiske programmer
i databasen.

En generel programoversigt bygges deterministisk fra de importerede sektioner
og krav. Mere specifikke spørgsmål kan besvares af Groq via
\code{get\_study\_plan}. Uanset tekstsvaret sendes studieplanen også som
strukturerede data, som browseren viser sektion for sektion. Ved fejl i
Groq/MCP genererer backend et statisk svar fra de samme databaserækker.

Sproget afgøres af den seneste forespørgsel. ``Computer Science study plan''
giver derfor engelske labels og links, mens ``Computer Science studieplan''
giver danske. De originale kursus- og sektionsnavne følger de importerede data,
men applikationens egne forklaringer oversættes.

\subsection{Specialiseringer}

Forespørgsler med danske, britiske eller amerikanske stavemåder af
``specialisering'' klassificeres særskilt. Systemet finder først programmet og
derefter eventuelt en navngiven specialisering. En oversigtsforespørgsel
returnerer alle importerede specialiseringer; en detaljeforespørgsel returnerer
krav, kursuspuljer og roller for den valgte specialisering.

Dette flow er primært deterministisk. Det reducerer risikoen for, at modellen
forveksler anbefalede kurser med obligatoriske kurser. Frontend anvender igen
\code{responseLanguage}, så blandt andet linkteksten ``View the specialization
at DTU'' vises på engelsk ved en engelsk forespørgsel.

\subsection{Åbne spørgsmål og opfølgninger}

Hvis ingen af de kendte intents passer, behandles beskeden som et åbent
spørgsmål. Groq kan stadig bruge de fire MCP-værktøjer, men systemprompten
kræver, at databasefakta hentes gennem et værktøj og ikke opfindes.

Opfølgninger bruger de tidligere brugerbeskeder som kontekst. Eksempelvis kan
brugeren først nævne ``Computer Science and Engineering'' og derefter spørge
``Hvilke specialiseringer har den?''. Begrænsningen er, at assistentens egne
svar ikke sendes med igen; nødvendig kontekst skal derfor kunne findes i
brugerens beskeder. Dette valg undgår samtidig, at meget lange kursuslister
overskrider inputgrænsen ved næste request.

\section{Svarformat, sprog og fejlhåndtering}

\code{ChatResponse} indeholder et kort tekstsvar, den forståede kontekst,
studieår og svar-sprog. Derudover kan det indeholde tre typer strukturerede
data: kursusanbefalinger, en studieplan og specialiseringer. Frontend behøver
derfor ikke parse modellens fritekst for at bygge kort og links.

Fejl håndteres på flere niveauer:

\begin{itemize}
  \item Ugyldigt input afvises af FastAPI med HTTP 422. Det gælder blandt andet
  en besked på mere end 800 tegn eller en request uden brugerbeskeder.
  \item Fejl i semantisk intent-klassifikation giver ikke i sig selv en fejl til
  brugeren; systemet fortsætter med den regelbaserede klassifikation.
  \item Fejl i OpenAI-embedding giver leksikalsk søgning som fallback.
  \item Fejl i Groq/MCP fører, hvor muligt, til en lokal kursussøgning eller et
  statisk databasebaseret studieplanssvar.
  \item Hvis hele HTTP-requesten fejler, viser den nuværende frontend en
  generel besked om, at der opstod en fejl. Status og detaljer logges kun i
  browserens konsol.
\end{itemize}

Det sidste punkt er en kendt designbegrænsning: brugeroplevelsen kan forbedres
ved at oversætte kendte HTTP-statusser til konkrete beskeder, eksempelvis at
en besked overskrider grænsen på 800 tegn.

\section{Drift og sikkerhedsgrænser}

Lokalt kan applikationen køres med Docker Compose. PostgreSQL-containeren
bruger pgvector, mens API-containeren først kører Alembic-migrationer og
derefter starter Uvicorn. I Vercel-miljøet kører FastAPI som en serverless
funktion, og databasekonfigurationen tager højde for en Supabase transaction
pooler.

Det almindelige REST-API er beskyttet med en delt \code{X-API-Key}, og MCP-
endpointet kræver et bearer-token, som sammenlignes sikkert. Chat-endpointet er
offentligt, så det kan bruges direkte fra browseren. Projektet har derimod ikke
applikationsstyret rate limiting på chatten. Det betyder, at misbrug primært
begrænses af hostingmiljøet og kvoterne hos Groq og OpenAI. For offentlig drift
bør der tilføjes rate limiting, overvågning og et budgetloft for eksterne
modelkald.

\section{Designvurdering}

Arkitekturen kombinerer deterministisk programlogik med sprogmodeller. Denne
hybrid er velegnet til domænet: databasen er sandhedskilden, mens modellen
bruges til at forstå fri tekst og formulere et læsbart svar. Strukturerede
studieplans- og specialiseringsdata begrænser risikoen for hallucinationer, og
fallbacks gør kursussøgningen mindre afhængig af eksterne tjenester.

De vigtigste kompromiser er flere netværkskald og dermed højere svartid, en
offentlig chat uden egen rate limiting samt en samtalehistorik, der kun består
af brugerbeskeder. Søgekvaliteten afhænger desuden af, at importerede data og
embeddings holdes opdaterede. Disse begrænsninger er tydelige og kan forbedres
uafhængigt, fordi præsentation, orkestrering, søgning og dataimport er adskilt.

\section{Konklusion}

DTU Course API er designet som en lagdelt, databasecentreret studievejleder.
Den seneste besked bestemmer brugerens aktuelle hensigt, mens tidligere
brugerbeskeder giver kontekst. Kursussøgning kombinerer fuldtekst og embeddings,
direkte opslag bruger kursusnummeret, og studieplaner samt specialiseringer
læses fra normaliserede kravmodeller. Groq fungerer som sproglig orkestrator
via afgrænsede MCP-værktøjer, men de officielle importerede DTU-data forbliver
grundlaget for svaret. Det giver et system, som både kan forstå naturligt sprog
og bevare en tydelig grænse mellem genereret tekst og faktuelle data.

\end{document}
